% Generated by code/interviews-1/03-speed-drills/python/gen_drill.py. Do not edit. French copy by draft/pilot-quant-fr/ivgen_fr.py.
\ivfamily{Additionner et soustraire des nombres entiers}{3, 8, 15, 16, 20, 31, 35, 41, 54, 55, 62, 72, 74, 75}{Travailler depuis la gauche, d'abord les centaines, et tenir un total courant.}
\ivfamily{Additionner et soustraire des décimaux}{7, 11, 39, 40, 42, 49, 57, 79}{Aligner la virgule mentalement, ou multiplier par 100 puis diviser par 100.}
\ivfamily{Multiplier par un chiffre}{14, 21, 36, 48, 51, 60, 63, 65, 73, 76}{Décomposer le plus grand nombre, ou arrondir puis soustraire quand c'est plus court~: $389 \times 7 = 400 \times 7 - 11 \times 7 = 2723$.}
\ivfamily{Nombres avec un raccourci}{1, 2, 5, 13, 29, 37, 50, 53, 59, 64, 67, 71}{$\times 25$~: multiplier par 100, puis diviser par 4. $\times 15$~: ajouter la moitié et multiplier par 10. $\times 11$~: multiplier par 10 et ajouter le nombre lui-même. $\times 99$~: multiplier par 100 et soustraire le nombre lui-même. $(n - d)(n + d) = n^2 - d^2$.}
\ivfamily{Division exacte}{6, 10, 27, 38, 44, 47, 66, 68, 70, 80}{$\div 5$~: doubler puis diviser par 10. $\div 25$~: multiplier par 4 et diviser par 100. $\div 12$ et $\div 15$~: diviser d'abord par 3.}
\ivfamily{Fractions}{12, 22, 24, 25, 26, 43, 45, 52, 61, 69}{Utiliser le plus petit dénominateur commun, multiplier en croix, simplifier avant de multiplier. Les réponses sont sous forme irréductible.}
\ivfamily{Pourcentages}{4, 9, 17, 19, 28, 30, 32, 46, 56, 58}{Partir de 10~\% et monter~: $15\%$ vaut $10\% + 5\%$, et $12.5\%$ vaut un huitième. $p\%$ de $n$ égale $n\%$ de $p$.}
\ivfamily{Produits de décimaux}{18, 23, 33, 34, 77, 78}{Multiplier les chiffres et compter les décimales. Traiter $0.25$, $0.75$, $1.5$ et $2.5$ comme des fractions~: $\tfrac14$, $\tfrac34$, $\tfrac32$, $\tfrac52$.}
